% 导言区
\documentclass[UTF8]{article} % book report letter
\usepackage{ctex}
\usepackage{verbatim}
\usepackage{amsmath}
\usepackage{graphicx}
\usepackage{palatino}
\usepackage{tikz}
\usepackage{float}

\title{\heiti 数学建模基础}
\author{\kaishu 庄梓伸}
\date{\kaishu \today}

% 正文区
\begin{document}
\maketitle

\section{层次分析法}

适用于评价类问题，其主要思路步骤如下
\begin{itemize}
    \item 明确评价目标
    \item 确定可选方案
    \item 查阅相关文献确定方案的评价指标
    \item 对各评价指标权重进行标度并得到判断矩阵
    \item 对判断矩阵进行一致性检验并适当修正
    \item 计算权重列写权重表进行分析
\end{itemize}

\subsection{标度及判断矩阵}

当面对多个评价指标时，我们往往无法综合地评判个指标的重要性，
因此分别比较各指标的权重，称为 标度 ，标度的含义如表1所示。

\begin{table}[ht]
    \centering
    \begin{tabular}{|c|c|}
        \hline
        标度 & 含义 \\
        \hline
        1 & 同等重要\\
        \hline
        3 & 稍微重要\\
        \hline
        5 & 比较重要\\
        \hline
        7 & 明显重要\\
        \hline
        9 & 极端重要\\
        \hline
        2、4、6、8 & 中值\\
        \hline
        倒数 & 相反\\
        \hline
    \end{tabular}
    \caption{标度的数值含义}
    \label{tab:table}
\end{table}

然后由两两指标的标度我们可以得到表2，该表写为矩阵形式，称为判断矩阵，形式如下。

\begin{equation}
    \frac{1}{2}
\end{equation}

我们称满足$a_{ij} a_{jk} = a_{ik}$的判断矩阵为一致矩阵。
往往，我们得到的判断矩阵不是一直矩阵，此时我们应该进行一致性检验。

\begin{table}[ht]
    \centering
    \begin{tabular}{|c|c|c|c|c|}
        \hline
        & 指标1 & 指标2 & 指标3 & ... \\
        \hline
        指标1 & 1 & $a/b$ & $a/c$ & ... \\
        \hline
        指标2 & $b/a$ & 1 & $b/c$ & ... \\
        \hline
        指标3 & $c/a$ & $c/b$ & 1 & ... \\
        \hline
        ... & ... & ... & ... & 1 \\
        \hline
    \end{tabular}
    \caption{标度表}
    \label{tab:table}
\end{table}

\subsection{一致性检验}



\subsection{权重表}

\begin{table}[H]
    \centering
    \begin{tabular}{|c|c|c|c|c|}
        \hline
        & 指标权重 & 方案1 & 方案2 & ... \\
        \hline
        指标1 & & & &\\
        \hline
        指标2 & & & &\\
        \hline
        ... & & & &\\
        \hline
    \end{tabular}
    \caption{权重表}
    \label{tab:table}
\end{table}

\end{document}
